\subsection[Super T-paths in a single fan triangulation]{Super $\boldsymbol{T}$-paths in a single fan triangulation}

Let $P$ be an $n$-gon and $T$ a fan triangulation. Let vertex $1$ be the fan center.

\begin{Lemma}[Schiffler] \label{lemma:Schiffler}
	For $2\leq i\leq n-1$, define
	\[\alpha_{i-1}:=(2,1,i,i+1,1,n\,|\,\dots),\qquad2\leqslant i\leqslant n-1.\]
	Then $\mathcal T^0_{2,n}=\{\alpha_i\colon 1\leq i\leq n-2\}$. See Example~{\rm \ref{exmp:T-Path-fan}}.
	
	Note that in case of $i=2$ or $i=n-1$, $\alpha_{i-1}$ collapses to a shorter $T$-path after removing backtracking. In the special case of $n=3$, there is a unique $T$-path in $\mathcal T^0_{2,n}$, i.e., $\alpha_2 = \alpha_{n-1}$, which consists of the single edge~$(2,3)$.
\end{Lemma}

\begin{Example}[ordinary $T$-paths of a fan]\label{exmp:T-Path-fan}\rm
	The following are the ordinary $T$-paths of a fan triangulation of an octagon. Notice that each of these $T$-paths surrounds (and is in bijection with) one of the triangles of $T$, as illustrated in yellow below:
\begin{center}
\input{example_T_path_in_fan}
\end{center}
\end{Example}

\begin{Lemma}\label{lm:super-path-single-fan}
	Every super $T$-path in $\mathcal T^1_{2,n}$ is of the following form
	\[(2,1,\theta_i,\theta_j,1,n\,|\,x_1,\sigma_i,\tau_{ij},\sigma_j,x_n)\]
	for $1\leq i<j\leq n-1$. See Example~{\rm \ref{exmp:super-T-path-fan}}.
\end{Lemma}

We illustrate this lemma with an example before proving it.

\begin{Example}\label{exmp:super-T-path-fan}\rm \allowdisplaybreaks
 The following are the ${4 \choose 2}=6$ different non-ordinary super $T$-paths (elements of $\mathcal{T}^1_{2,6}$) of a single-fan hexagon:
\begin{center}
\input{example_super_T_path_in_fan}
\end{center}
In this example, $\lambda_{26}$ can be expressed as
 \begin{gather*}
 \lambda_{26} = \frac{\lambda_{23}\lambda_{16}}{\lambda_{31}}
 + \frac{\lambda_{21}\lambda_{34}\lambda_{16}}{\lambda_{13}\lambda_{41}}
 + \frac{\lambda_{21}\lambda_{45}\lambda_{16}}{\lambda_{14}\lambda_{51}}
 + \frac{\lambda_{21}\lambda_{56}}{\lambda_{15}} \\
 \phantom{\lambda_{26}=} + \lambda_{21} \sqrt{\frac{\lambda_{23}}{\lambda_{12}\lambda_{13}}} \, \boxed{123} \, \sqrt{\frac{\lambda_{34}}{\lambda_{13}\lambda_{14}}} \, \boxed{134} \, \lambda_{16}
+ \lambda_{21} \sqrt{\frac{\lambda_{34}}{\lambda_{13}\lambda_{14}}} \, \boxed{134} \, \sqrt{\frac{\lambda_{45}}{\lambda_{14}\lambda_{15}}} \, \boxed{145} \, \lambda_{16} \\
 \phantom{\lambda_{26}=} + \lambda_{21} \sqrt{\frac{\lambda_{45}}{\lambda_{14}\lambda_{15}}} \, \boxed{145} \, \sqrt{\frac{\lambda_{56}}{\lambda_{15}\lambda_{16}}} \, \boxed{156} \, \lambda_{16}
+ \lambda_{21} \sqrt{\frac{\lambda_{23}}{\lambda_{12}\lambda_{13}}} \, \boxed{123} \, \sqrt{\frac{\lambda_{45}}{\lambda_{14}\lambda_{15}}} \, \boxed{145} \, \lambda_{16} \\
 \phantom{\lambda_{26}=} + \lambda_{21} \sqrt{\frac{\lambda_{34}}{\lambda_{13}\lambda_{14}}} \; \boxed{134} \; \sqrt{\frac{\lambda_{56}}{\lambda_{15}\lambda_{16}}} \, \boxed{156} , \lambda_{16}
+ \lambda_{21} \sqrt{\frac{\lambda_{23}}{\lambda_{12}\lambda_{13}}} \, \boxed{123} \, \sqrt{\frac{\lambda_{56}}{\lambda_{15}\lambda_{16}}} \, \boxed{156} \, \lambda_{16}.
 \end{gather*}
 The first four terms are the weights of all ordinary $T$-paths (as described in Example~\ref{exmp:T-Path-fan}),
 and the remaining six terms correspond to the $T$-paths pictured above (in the same order).
 It can also be written as
 \begin{gather*}
 \lambda_{26} = \frac{\lambda_{23}\lambda_{16}}{\lambda_{31}}
 + \frac{\lambda_{21}\lambda_{34}\lambda_{16}}{\lambda_{13}\lambda_{41}}
 + \frac{\lambda_{21}\lambda_{45}\lambda_{16}}{\lambda_{14}\lambda_{51}}
 + \frac{\lambda_{21}\lambda_{56}}{\lambda_{15}}
 + \lambda_{16}\lambda_{21} \, \widetilde{\boxed{123}} \; \widetilde{\boxed{134}}
 + \lambda_{16}\lambda_{21} \, \widetilde{\boxed{134}} \; \widetilde{\boxed{145} }\\
 \hphantom{\lambda_{26} =}{}
 + \lambda_{21}\lambda_{16}\, \widetilde{\boxed{145}} \; \widetilde{\boxed{156} }
 + \lambda_{16}\lambda_{21} \, \widetilde{\boxed{123}} \; \widetilde{ \boxed{145}}
 + \lambda_{16}\lambda_{21} \, \widetilde{\boxed{134}}\;\widetilde{\boxed{156}}
 + \lambda_{16}\lambda_{21}\,
 \widetilde{ \boxed{123} }\, \widetilde{ \boxed{156}}\,,
 \end{gather*}
 where the $\widetilde{\boxed{ijk}}$ are the weights of the corresponding $\sigma$-edges (following the notation of Corollary~\ref{cor:laurent}).
\end{Example}

\begin{proof}[Proof of Lemma~\ref{lm:super-path-single-fan}]
 A $\mathcal T^1_{2,n}$-path must use one of the $\sigma$-edges, therefore the first step in the super $T$-path needs to be $(2,1)$.
 Then after traveling through $\sigma_i$, what follows must be $\tau_{ij}$ which leads the $T$-path to another internal vertex $\theta_j$.
 The next step is an even step hence has to be a $\sigma$-edge which will take us to vertex $1$:
 \[(2,1,\theta_i,\theta_j,1,\dots\,|\,x_1,\sigma_i,\tau_{ij},\sigma_j,\dots).\]
 Now we are at vertex $1$ and have completed even number of steps, therefore the rest of the this super $T$-path must be a super $T$-path
 from~$1$ to~$n$~-- clearly there is only one possibility which is the single edge $(1,n)$.
 Hence a super $T$-path in $\mathcal T^1_{2,n}$ must have the form
 \begin{gather*}
 (2,1,\theta_i,\theta_j,1,\dots\,|\,x_1,\sigma_i,\tau_{ij},\sigma_j,\dots)+(\dots,1,n\,|\,\dots,x_n)\\
 \qquad{} =(2,1,\theta_i,\theta_j,1,n\,|\,x_1,\sigma_i,\tau_{ij},\sigma_j,x_n).\tag*{\qed}
 \end{gather*}\renewcommand{\qed}{}
\end{proof}

\section{Proof of Theorem~\ref{thm:main}}
\label{sec:proof}

In this section we prove our main theorem.
It turns out that the default orientation guarantees a positive sign on all terms of the expansion of
$\lambda$-lengths (Theorem~\ref{thm:main}), and is also preserved by induction.

Our proof has three parts: we first prove the case of single fan triangulations, and then prove the case of zig-zag triangulations. Finally we prove Theorem~\ref{thm:main} in full generality by combining the two cases mentioned above.

Before proving our theorem, we state the following results that will be used in our proofs.

\begin{Proposition} \label{lemma:theta32}
 Let $A$, $\beta,$ and $\Sigma$ be elements in the super algebra $\mathcal A$, which for convenience, we assume is written as
 $\mathcal A = {\mathbb R}\bigl[x_1^{\pm\frac{1}{2}},\dots,x_{2n+3}^{\pm\frac{1}{2}} \,|\, \theta_1,\dots,\theta_{n+1}\bigr]$ as in Definition~{\rm \ref{def:weight}}.
 Further, we assume that $A$ is an even element with a non-zero body,\footnote{The \emph{body} of an element $A$ of super algebra $\mathcal A$ is the constant term when expanded out in terms of the $\theta_i$'s.}
 and that both $\beta$ and $\Sigma$ are odd elements of~$\mathcal A$. Then we have
 \[\sqrt{A+\beta\Sigma}= \sqrt{A}+\frac{\beta}{2\sqrt{A}}\Sigma,\]
 where the square root of $A$ is taken to be the positive square root.
\end{Proposition}

\begin{Remark}\rm
 Since $A$ is an even element, its positive square root is well-defined as the choice such that the body of $\sqrt{A}$ is the positive square root of the body of $A$.
\end{Remark}

\begin{proof}
 Squaring the right-hand side:
 \[\left(\sqrt{A}+\frac{\beta}{2\sqrt{A}}\Sigma\right)^2=A+2\sqrt{A}\cdot \frac{\beta}{2\sqrt{A}}\cdot \Sigma +\left(\frac{\beta}{2\sqrt{A}}\Sigma\right)^2=A+\beta\Sigma.\]
 This clearly equals the square of the left-hand side.
\end{proof}

Using Lemma~\ref{lemma:theta32}, we can rewrite the super Ptolemy relations in a more symmetrical form.
\begin{Proposition}
The super Ptolemy relations described in Figure~{\rm \ref{fig:super_ptolemy}} can be written as follows
\begin{gather}
\theta'\sqrt{ef} =	\theta\sqrt{bd}+\sigma \sqrt{ac},\label{eq:5}\\
\sigma'\sqrt{ef} =\sigma\sqrt{bd}-\theta\sqrt{ac}.\label{eq:6}
\end{gather}
\end{Proposition}
\begin{proof}We have
\begin{gather*}
\theta'\sqrt{ef} =\theta'\sqrt{ab+cd+\sqrt{abcd}\sigma\theta}
 =\theta'\sqrt{ab+cd+\sqrt{abcd}\sigma'\theta'}
 =\theta'\sqrt{ab+cd}
\end{gather*}
and
\begin{gather*}
	\theta' = {{ \theta\sqrt{bd}+\sigma \sqrt{ac}}\over{\sqrt{ac+bd}}},\qquad
	\theta'\sqrt{ac+bd} = \theta\sqrt{bd}+\sigma \sqrt{ac}.
\end{gather*}
	
Putting these two equations together gives equation~\eqref{eq:5}:
\begin{align*}
\theta'\sqrt{ef}= \theta\sqrt{bd}+\sigma \sqrt{ac}.
\end{align*}
Equation~\eqref{eq:6} can be derived in a similar way.
\end{proof}

\subsection{Proof of Theorem~\ref{thm:main} for a single fan}

For sake of readability, in the below, we use $\boxed{ijk}$ to denote the $\mu$-invariant associated to the triple of
ideal points $(i,j,k)$, subject to Remark~\ref{rem:mu-invariants}, while the $\lambda$-length of a pair $(i,j)$ will be
denoted $\lambda_{ij}$. We will also sometimes use $\boxed{ijk}$ to denote the internal vertex of the auxiliary graph associated to $(i,j,k)$, when talking about super $T$-paths.

First, we will prove the main theorem in the case of a single fan triangulation.

\begin{figure}[h] \centering
 \input{proof_aux_graph}
 \caption{Proof of Theorem~\ref{thm:mu_single_fan}.} \label{fig:mu_single_fan}
\end{figure}

\begin{Theorem} \label{thm:mu_single_fan}
 Consider a single fan triangulation of an $n$-gon as depicted in Figure~{\rm \ref{fig:mu_single_fan}}. We continue our convention of using the default orientation,
 which for a fan triangulation means that arrows on all internal diagonals point away from the fan center.
 After performing the super Ptolemy relations for the flips of arcs $(1,3)$, $(1,4)$, $\dots$, $(1,k-1)$, we have
 \begin{enumerate}\itemsep=0pt
 \item[$(a)$] $\displaystyle \sqrt{\frac{\l2k}{\l12\l1k}} \; \boxed{12k} = \sum_{i=1}^{k-2}\wt(\sigma_i)$,
 \item[$(b)$] $\displaystyle \lambda_{2k} = \sum_{t \in \mathcal{T}_{2,k}} {\rm wt}(t)$,
 \end{enumerate}
 where ${\rm wt}(\sigma_i)$ and ${\rm wt}(t)$ are defined in Definition~{\rm \ref{def:weight}}, and as defined in Definition~{\rm \ref{def:super-T}}, $\mathcal{T}_{2,k}$ denotes the set of super $T$-paths from vertex $2$ to vertex~$k$.
\end{Theorem}

\begin{proof}
 We will induct on $k$. We begin with the case of $k=3$ where statements (a) and (b) follow immediately since
 ${\rm wt}(\sigma_1) = \sqrt{\frac{\lambda_{23}}{\lambda_{12}\lambda_{13}}} \, \boxed{123}$\,, and the unique super $T$-path from vertex $2$ to vertex $3$
 is simply the ordinary $T$-path $t = (2,3)$.

 In general, for $k\geq 4$, after flipping arcs $(1,3),\dots,(1,k-2)$, the next flip of arc $(1,k-1)$ will be inside the following quadrilateral:
 \begin{center}
 \begin{tikzpicture}
 \draw (1,0) -- (0,1) -- (-1,0) -- (0,-1) -- cycle;
 \draw (-1,0) --node {\midarrow} (1,0);

 \draw (-1,0) node[left] {$1$};
 \draw (1,0) node[right] {$k-1$};
 \draw (0,-1) node[below] {$2$};
 \draw (0,1) node[above] {$k$};
 \end{tikzpicture}
 \end{center}
 Note that the $\mu$-invariant $\boxed{1,k-1,k}$ is in the initial triangulation, but $\boxed{1,2,k-1}$ is not\footnote{Except for the special case of $k=4$.}.
 However, we assume by induction that $\boxed{1,2,k-1}$ is given by part $(a)$.

 First we prove part (b).
 The super Ptolemy relation (equation~\eqref{eq:mutation}) says that $\lambda_{2k}$ is given by
 \[
 \lambda_{1,k-1}\lambda_{2k} = \lambda_{12}\lambda_{k-1,k} + \lambda_{1k}\lambda_{2,k-1}
 + \sqrt{\lambda_{12}\lambda_{2,k-1}\lambda_{k-1,k}\lambda_{1k}} \, \boxed{1,2,k-1} \, \boxed{1,k-1,k}\,.
 \]
 After dividing by $\lambda_{1,k-1}$, we get a formula for $\lambda_{2,k}$. On the right-hand side, the first term gives
 $\lambda_{21}\lambda_{1,k-1}^{-1}\lambda_{k-1,k}$, which is clearly an ordinary $T$-path. The second term, by induction, is
 \[ \sum_{t \in \mathcal{T}_{2,k-1}} {\rm wt}(t) \lambda_{k-1,1}^{-1} \lambda_{1k}. \]
 Taking the ordinary $T$-paths in $\mathcal{T}^0_{2,k-1}$, and appending the arcs $(k-1,1)$ and $(1,k)$ give the rest of the ordinary $T$-paths
 in $\mathcal{T}^0_{2k}$. (See Lemma~\ref{lemma:Schiffler} and note that appending arc $(k-1,1)$ as an even step may in fact yield a backtrack that cancels out the final step of an ordinary $T$-path in~$\mathcal{T}^0_{2k}$.) The terms coming from $\mathcal{T}^1_{2,k-1}$, multiplied by~$\frac{\lambda_{1k}}{\lambda_{1,k-1}}$, similarly give the super $T$-paths
 in~$\mathcal{T}^1_{2k}$ which do not involve the triangle~$(1,k-1,k)$ (see Lemma~\ref{lm:super-path-single-fan}).

 The last term, using part (a) and induction on $\boxed{1,2,k-1}$\,, is
 \begin{align*}
 \frac{\sqrt{\lambda_{12}\lambda_{2,k-1}\lambda_{k-1,k}\lambda_{1k}}}{\lambda_{1,k-1}} \, \boxed{1,2,k-1} \, \boxed{1,k-1,k}
 &= \lambda_{21} \sum_{i=1}^{k-3} {\rm wt}(\sigma_i) {\rm wt}(\sigma_{k-2}) \lambda_{1k}.
 \end{align*}
 These are the weights of the super $T$-paths which use the last triangle, namely $(1,k-1,k)$ (again, see Lemma~\ref{lm:super-path-single-fan}).

 Note that the product ${\rm wt}(\sigma_i) {\rm wt}(\sigma_{k-2})$ is in the positive ordering,
 since the two $\mu$-invariants appear in the counter-clockwise order around the fan center.

 Now, we examine part (a). Looking at the same quadrilateral as above, the super Ptolemy relation for $\mu$-invariants (equation~\eqref{eq:5}) says that
 \[ \boxed{12k} \, \sqrt{\lambda_{1,k-1}\lambda_{2k}} = \boxed{1,k-1,k} \, \sqrt{\lambda_{12}\lambda_{k-1,k}} + \boxed{1,2,k-1} \, \sqrt{\lambda_{1k}\lambda_{2,k-1}}. \]

 Dividing by $\sqrt{\l12\lambda_{1,k-1}\l1k}$, commuting the $\lambda$-lengths past the $\mu$-invariants, we get
 \[\sqrt{\l2k\over \l12\l1k }\, \boxed{12k}=\sqrt{\l k{,k-1} \over \l1{,k-1}\l1k }\,\boxed{1,k-1,k} + \sqrt{\l2{,k-1} \over\l12\l1{,k-1}}\,\boxed{1,2,k-1}\,.\]
 The first term on the right-hand side is simply ${\rm wt}(\sigma_{k-2})$. By induction, the second term is $\sum_{i=1}^{k-3} {\rm wt}(\sigma_i)$. Therefore we have
 \[\sqrt{\l2k\over \l12\l1k }\, \boxed{12k}=\wt(\sigma_{k-2}) + \left(\sum_{i=1}^{k-3}\wt(\sigma_i)\right)=\sum_{i=1}^{k-2}\wt(\sigma_{i})\]
 as desired.
\end{proof}

\subsection{Proof of Theorem~\ref{thm:main} for a zig-zag triangulation}

Next, we prove our main theorem for the case of a zig-zag triangulation.

\begin{Theorem} \label{thm:proof_zigzag}
 Consider a zigzag triangulation $T$ of an $n$-gon as depicted in Figure~{\rm \ref{fig:proof_zigzag}}. We consider all the vertices
 except for $1$ and $n$ to be fan centers, so that $c_{i}$ is labelled $i+1$ for $1\leq i\leq n-2$.

 After flipping the arcs $(n-1,n-2),(n-2,n-3),\dots,(k+2,k+1)$, we have
 \begin{enumerate}\itemsep=0pt
 \item[$(a)$] $\displaystyle \sqrt{\lambda_{kn}\lambda_{k+1,n}\over\lambda_{k,k+1}} \, \boxed{k,k+1,n} = \sum_{i=k}^{n-2} \sum_{t\in\mathcal T_{n,i+1}} \wt(t) \wt(\sigma_i)$,
 \item[$(b)$] $\displaystyle \lambda_{k,n} = \sum_{t \in \mathcal{T}_{kn}} {\rm wt}(t)$.
 \end{enumerate}
\end{Theorem}

\begin{figure}[h]\centering
\input{figure_proof_zigzag}
\caption{Proof of Theorem~\ref{thm:proof_zigzag}.
{Left:} zig-zag triangulation with the default orientation.
{Right:} the corresponding auxiliary graph.}\label{fig:proof_zigzag}
\end{figure}

\begin{proof}
 We assume that vertices $n$, $n-1$, and $n-2$ are oriented as in Figure~\ref{fig:proof_zigzag}. The case that they are oriented oppositely is similar.

 We will induct (backwards) on $k$.
 The base case of the induction is the case of a single triangle, when $k=n-2$.
 Since the edge $(n-2,n)$ is already in the triangulation, $\lambda_{n-2,n}$ is already one of the generators. Clearly the only $T$-path
 in this case is the single edge $(n-2,n)$ (when zero flips have been performed). This establishes part (b) for the base case.
 For part (a), the left-hand side is
 \[ \sqrt{\frac{\lambda_{n-2,n}\lambda_{n-1,n}}{\lambda_{n-1,n-2}}} \, \boxed{n-2,n-1,n}\,. \]
 For the right-hand side, there is only a single term in this sum. This is because $i$ only takes the value $n-2$,
 and the only $T$-path from $n-1$ to $n$ is the single edge $(n-1,n)$. Thus the right-hand side is
 \begin{align*}
 \lambda_{n-1,n}   \wt(\sigma_{n-2}) &= \lambda_{n-1,n}   \sqrt{\frac{\lambda_{n-2,n}}{\lambda_{n-2,n-1}\lambda_{n-1,n}}} \, \boxed{n-2,n-1,n} \\
 &= \sqrt{\frac{\lambda_{n-2,n}\lambda_{n-1,n}}{\lambda_{n-2,n-1}}} \, \boxed{n-2,n-1,n}\,.
 \end{align*}
 This establishes part (a) for the base case. We now assume that $1 \leq k \leq n-3$.

 After flipping the arcs $(n-1,n-2)$, $(n-2,n-3), \dots, (k+3,k+2)$,\footnote{For the case of $k=n-3$, no arcs have yet been flipped.} we will have one of the two following quadrilaterals, depending on the parity of $n-k$:
 \begin {center}
 \begin {tikzpicture}
 \draw (1,0) -- (0,1) -- (-1,0) -- (0,-1) -- cycle;
 \draw (-1,0) --node {\midarrow} (1,0);
 \draw (-1,0) node[left] {$k+1$};
 \draw (1,0) node[right] {$k+2$};
 \draw (0,-1) node[below] {$n$};
 \draw (0,1) node[above] {$k$};
 \draw (3,0) node {or};
 \begin {scope}[shift={(6,0)}]
 \draw (1,0) -- (0,1) -- (-1,0) -- (0,-1) -- cycle;
 \draw (1,0) --node {\reflectbox{\midarrow}} (-1,0);
 \draw (-1,0) node[left] {$k+2$};
 \draw (1,0) node[right] {$k+1$};
 \draw (0,-1) node[below] {$n$};
 \draw (0,1) node[above] {$k$};
 \end {scope}
 \end {tikzpicture}
 \end {center}


 Now we flip the edge $(k+2, k+1)$ while applying the Ptolemy relation equation~\eqref{eq:5}. In both cases pictured above, the triangle $(k,k+1,n)$ will play
 the role of $\theta'$ (it will be on the left, looking in the direction of the arrow after the flip).
 And because of the opposite orientations, application of equation~\eqref{eq:5} in both pictures gives
 \begin{gather*}
 \sqrt{\lambda_{kn} \lambda_{k+1,k+2}} \, \boxed{k,k+1,n} \\
 \qquad{} = \sqrt{\lambda_{k+1,n} \lambda_{k,k+2}} \, \boxed{k,k+1,k+2} + \sqrt{\lambda_{k,k+1} \lambda_{k+2,n}} \, \boxed{k+1,k+2,n}\,.
 \end{gather*}


 Multiplying both sides by $\sqrt{\frac{\lambda_{k+1,n}}{\lambda_{k,k+1}\lambda_{k+1,k+2}}}$, we get
 \begin{gather*}
 \begin{split}
 & \sqrt{\frac{\lambda_{kn} \lambda_{k+1,n}}{\lambda_{k,k+1}}} \, \boxed{k,k+1,n}\\
& \qquad{}
 = \lambda_{k+1,n} \sqrt{\frac{\lambda_{k,k+2}}{\lambda_{k,k+1}\lambda_{k+1,k+2}}} \, \boxed{k,k+1,k+2}
 + \sqrt{\frac{\lambda_{k+1,n} \lambda_{k+2,n}}{\lambda_{k+1,k+2}}} \, \boxed{k+1,k+2,n}\,.
 \end{split}
 \end{gather*}


 First we examine the first term on the right-hand side. By induction, $\lambda_{k+1,n}$ is the weighted sum of super $T$-paths from $k+1$ to $n$.
 Notice that $\sqrt{\frac{\lambda_{k,k+2}}{\lambda_{k,k+1} \lambda_{k+1,k+2}}} \, \boxed{k,k+1,k+2}$ is the weight of the $\sigma$-step $\sigma_k$, going from
 vertex $k+1$ into the triangle labelled $\theta_k$.
 Therefore this first term is equal to
 \begin{equation} \label{eq:zigzag_1st_term}
 \sum_{t \in \mathcal{T}_{k+1,n}} \wt(t) \wt(\sigma_k)\,.
 \end{equation}

Next, we focus on the second term of the right-hand side: $\sqrt{\frac{\lambda_{k+1,n} \lambda_{k+2,n}}{\lambda_{k+1,k+2}}} \, \boxed{k+1,k+2,n}\,$.
 By induction, this is equal to
 \begin{equation*} %\label{eq:zigzag_2nd_term}
 \sqrt{\frac{\lambda_{k+1,n} \lambda_{k+2,n}}{\lambda_{k+1,k+2}}} \, \boxed{k+1,k+2,n} = \sum_{i=k+1}^{n-2} \sum_{t \in \mathcal T_{n,i+1}} \wt(t) \wt(\sigma_i).
 \end{equation*}
 Adding the terms from equation~\eqref{eq:zigzag_1st_term} gives the result for part~(a).

 We should verify that the $\mu$-invariants in each product $\wt(t)\wt(\sigma_k)$ are in the correct positive ordering. {We use the same viewpoint as in Remark~\ref{rem:pos_order}.} By induction, the $\mu$-invariants
 in each term of $\wt(t)$ occur in the positive ordering of the smaller polygon which contains triangles $\theta_{k+1},\dots,\theta_{n-2}$.
 By construction, the positive ordering of the slightly larger polygon (which additionally contains $\theta_k$) is obtained from the smaller by
 placing $\theta_k$ either at the beginning or end of the previous order. But since $\wt(t)$ is an even element, it is central, and so
 $\wt(t)\wt(\sigma_k) = \wt(\sigma_k)\wt(t)$, and we may choose whichever one gives the correct positive ordering.

 Now we will prove part (b). Again after flipping the arcs $(n-1,n-2), (n-2,n-3), \dots$, $(k+3,k+2)$, we are in one of the
 two pictures from before. For the picture on the left, an application of equation~\eqref{eq:mutation}, followed by division on both sides by $\lambda_{k+1,k+2}$, gives
 \begin{gather*}
 \l kn =
 \underbrace{
 \frac{\lambda_{k,k+2} \lambda_{k+1,n}}{\lambda_{k+1,k+2}}}_{\text{part 1}}
 +
 \underbrace{\frac{\lambda_{k,k+1} \lambda_{k+2,n}}{\lambda_{k+1,k+2}}}_{\text{part 2}}\nonumber \\
 \hphantom{\l kn =}{}+
 \underbrace{\sqrt{\frac{\lambda_{k+1,n} \lambda_{k+2,n}}{\lambda_{k+1,k+2}}}\; \boxed{k+1,k+2,n}}_{\text{part 3}}
 \cdot
 \underbrace{\sqrt{\frac{\lambda_{k,k+1} \lambda_{k,k+2}}{\lambda_{k+1,k+2}}}\; \boxed{k,k+1,k+2}}_{\text{part 4}}.%\label{eq:proof_zigzag_lambda}
 \end{gather*}
 Note that $\boxed{k+1,k+2,n}$ lies on the right of the oriented diagonal, and hence plays the role of~$\sigma$. Analogously,
 $\boxed{k,k+1,k+2}$ lies on the left and plays the role of $\theta$. Applying equation~\eqref{eq:mutation} to the picture on the right gives the same thing, except that parts~(3) and~(4) appear
 in the opposite order.

 First we will explain that the sum of parts~(1) and~(2) is the weighted sum of all super $T$-paths from $k$ to $n$
 that do not contain a $\tau$-edge which crosses $(k+1,k+2)$.

 Suppose a path $t \in \mathcal{T}_{kn}$ does not contain a $\tau$ edge crossing $(k+1,k+2)$ (i.e., does not have a super step starting with $\theta_k$).
 Then the first step of $t$ is either edge $(k,k+1)$ or $(k,k+2)$, and the remainder of $t$ lies in the smaller polygon below the diagonal $(k+1,k+2)$.
 There are two cases.

 The second edge might be $(k+1,k+2)$, in which case the remainder of $t$ (after removing the first two edges) is a super $T$-path
 in either $\mathcal{T}_{k+1,n}$ or $\mathcal{T}_{k+2,n}$. This is depicted in the left of Figure~\ref{fig:zigzag_cases}.
 Clearly all of these paths occur as terms in parts~(1) and~(2).

 However, there are more terms in parts (1) and~(2);
 namely those in which the denominator $\lambda_{k+1,k+2}$ cancels a contribution in the numerator. This is the
 second case, in which the second step in $t$ is \emph{not} the edge $(k+1,k+2)$. This is depicted in the right of Figure~\ref{fig:zigzag_cases}.
 In this case, replacing the first edge of~$t$ with $(k+1,k+2)$ gives
 a super $T$-path in either $\mathcal{T}_{k+1,n}$ or $\mathcal{T}_{k+2,n}$, and these are the remaining terms in parts~(1) and~(2) in which the denominator cancels.

 \begin{figure}[h]\centering
 \input{figure_zigzag_cases}
 \caption{{Left:} Removing the first two edges gives a path in $\mathcal{T}_{k+2,n}$.
 {Right:} Replacing the first edge $(k,k+2)$ with $(k+1,k+2)$ gives a path in $\mathcal{T}_{k+1,n}$.} \label{fig:zigzag_cases}
 \end{figure}

 Now we examine parts (3) and~(4). By induction, part (3) is given by the formula in part~(a):
 \[
 \sqrt{\frac{\lambda_{k+1,n} \lambda_{k+2,n}}{\lambda_{k+1,k+2}}}\, \boxed{k+1,k+2,n}
 = \sum_{i=k+1}^{n-2} \sum_{t \in \mathcal{T}_{i+1,n}} \wt(t) \wt(\sigma_i).
 \]
 Part (4) is equal to $\lambda_{k,k+1} \wt(\sigma_k)$, which is the weight of the first two steps of any super $T$-path $t \in \mathcal{T}_{kn}$ which \emph{does}
 contain a $\tau$-step crossing $(k+1,k+2)$. The terms of part~(3) are all possible ways to complete such a super $T$-path (after joining them by
 the appropriate $\tau$-step, which has weight~1).

\looseness=-1 Again, as in part~(a), we must check that these expressions are written in the correct positive ordering. As we noted above
 in the discussion of part~(a), the positive ordering of the smaller polygon (below the diagonal $(k+1,k+2)$) agrees with
 the positive ordering in the larger polygon. Therefore any factors appearing the formulas obtained above can be assumed to be in
 the correct positive ordering. In parts~(1) and~(2), two of the three factors are single edges in the triangulation, and the terms of the third factor
 (either $\lambda_{k+1,n}$ or $\lambda_{k+2,n}$) appear in positive order by induction. As was discussed in part $(a)$, the terms of
 part~(3) can be written as either $\wt(t)\wt(\sigma_i)$ or $\wt(\sigma_i)\wt(t)$, whichever is correct. Part~(4) only contains a single $\mu$-invariant.
 So all that needs to be checked is that parts~(3) and~(4) occur in the correct order, depending on whether $\theta_k$ comes first or last
 in the positive ordering. But this is precisely the difference between the left and right pictures above (depending on whether vertex~$k$ or~$n$ is on the left of the oriented edge $k+1 \to k+2$), and parts~(3) and~(4) are positioned differently in the two cases.
\end{proof}

\subsection{Proof of Theorem~\ref{thm:main} for generic triangulations}
By a generic triangulation, we mean a polygon in which every triangle has at least one boundary edge. This is the same hypothesis that appears in Proposition~\ref{prop:equiv_rel} because of Remark~\ref{rem:subtriang}. Given a generic triangulation $T$ with $n$ fans, we first apply the flip sequence in Theorem~\ref{thm:mu_single_fan} on each of the fans,
which result in a zig-zag (sub)triangulation $T'$ whose vertices are the fan centers of~$T$ (including $c_0$ and $c_{N+1}$). See Figure~\ref{fig:main_flip_sequence}.

\begin{figure}[h]
\centering
\input{figure_main_flip_sequence.tex}
\caption{Flipping edges in each fan, to turn a generic triangulation $T$ into a zig-zag triangulation $T'$.}\label{fig:main_flip_sequence}
\end{figure}

We will then derive our ultimate formula for $\lambda_{ab}$ in $T$ via a combination of Theorem~\ref{thm:mu_single_fan} and Theorem~\ref{thm:proof_zigzag}.
Using Theorem~\ref{thm:proof_zigzag}, we can express $\lambda_{ab}$ in terms of super $T$-paths on $T'$.
This expression uses certain $\lambda$-lengths and $\mu$-invariants that are not in $T$, so we will substitute their expansion
in terms of $T$ using Theorem~\ref{thm:mu_single_fan}.

\begin{figure}[h]
\input{figure_main_flip_sequence_labelled.tex}
\caption{The auxiliary graphs for $T$ and $T'$, with $\tau$-edges omitted.}\label{fig:main_flip_sequence_labelled}
\end{figure}

We consider the fans of $T$ as subtriangulations, and denote them as $F_1,\dots,F_N$.
We denote the $\mu$-invariants of $T'$ by $\theta_1',\dots,\theta_N'$ and their corresponding $\sigma$-edges to be $\sigma_1',\dots,\sigma_N'$.
We denote the $\mu$-invariants of the $i$-th fan of $T$ by $\theta_i^1,\theta_i^2,\dots$ (ordered counterclockwise around the fan center),
and the corresponding $\sigma$-edges $\sigma_i^1,\sigma_i^2,\dots$. See Figure~\ref{fig:main_flip_sequence_labelled}.

When we substitute super $T$-path expressions for the fans in $T$ into $T'$, we only need to consider two cases:
(1) substitute a boundary edge $(c_{i-1},c_{i+1})$ for $1\leqslant i\leqslant n$,
and (2) substitute a super-step $(\dots,\sigma_i',\tau_{ij},\sigma_j',\dots\,|\,\dots)$,
because every super $T$-path is a concatenation of super steps and complete ordinary $T$-paths.

\begin{enumerate} \itemsep=0pt
 \item[(1)] Suppose we were to replace $\lambda_{c_{i-1},c_{i+1}}$ with its super $T$-path expansion in the $i$-th fan.
 In the super $T$-path of $T'$, the edge $(c_{i-1},c_{i+1})$ must be an odd step because it does not cross the arc $(a,b)$.
 Therefore when we replace it with a super $T$-path in~$F_i$, the indexing agrees up to parity.
Moreover, if an edge crosses the arc $(c_{i-1},c_{i+1})$ in $F_i$, then it must also cross the arc $(a,b)$ in the bigger triangulation~$T$.
 Therefore the axiom (T5) is satisfied, and it is straightforward to verify that all other axioms will follow.

 \item[(2)] We observe that the left-hand side of Theorem~\ref{thm:mu_single_fan}(a)
 equals $\wt(\sigma_i^\prime)$ with respect to $T^\prime$. Using the rest of Theorem~\ref{thm:mu_single_fan}(a), we get the equality
 \[\wt(\sigma_i^\prime)=\sum_{j}\wt\big(\sigma_i^j\big),\]
 i.e., we have that the weight of $\sigma_i^\prime$ in $T^\prime$ is equal to the weighted sum of all $\sigma$-edges in the fan~$F_i$.
 This means that a super step $(\dots,\sigma_i^\prime,\tau_{ij},\sigma_j^\prime,\dots)$ will be expanded into the sum of all
 super steps from a face of $F_i$ to a face of $F_j$. This clearly preserves all axioms of super $T$-paths.
\end{enumerate}

For the converse, we need to prove that every super $T$-path in $T$ can be obtained by such a~substitution.
First, this is clearly true for ordinary $T$-paths, or an ordinary sub-path of a super $T$-path. Therefore we only need to consider the super steps.
Suppose we have a super-step using two $\sigma$-steps $\sigma^*$ and $\sigma^\bullet$, if $\sigma^*$ and $\sigma^\bullet$ are in the same fan,
say $F_i$, then the super step is part of the expansion of $\lambda$-length of $(c_{i-1},c_{i+1})$.
If the $\sigma^*\in F_i$ and $\sigma^\bullet\in F_j$ are in different fans, then the super step came from the super-step
$(\dots,\sigma_i^\prime,\tau_{ij},\sigma_j^\prime,\dots)$ in $T'$.

